\section{估值方法不能一刀切}

两家银行采用PB--ROE框架；徐工机械采用2026E PE并用现金流校验；中兴采用合并PE但只给算力期权信用；江波龙采用峰值周期PE并显式折价；恒瑞采用SOTP；工业富联采用预告后合并PE并用AI服务器、800G交换机和H2利润校验。七只核心若统一使用PE，会同时制造银行、周期股、低毛利制造和创新药的估值错配。

\begin{exhibitbox}[Exhibit 22：双篮子最终估值表]
\begin{tabularx}{\textwidth}{L{2.0cm}R{1.35cm}R{1.35cm}R{1.35cm}R{1.45cm}R{1.45cm}R{1.45cm}X}
\toprule
\textbf{公司} & \textbf{现价} & \textbf{熊} & \textbf{基准} & \textbf{牛} & \textbf{最终目标} & \textbf{空间} & \textbf{动作} \\
\midrule
渝农商行 & 6.30 & 6.39 & 8.18 & 9.59 & 8.01 & +27.1\% & 核心观察/回撤配置 \\
徐工机械 & 8.43 & 8.52 & 11.36 & 13.49 & 11.08 & +31.4\% & 回撤配置/业绩验证 \\
沪农商行 & 7.66 & 7.76 & 9.88 & 10.73 & 9.69 & +26.5\% & 收益型回撤配置 \\
中兴通讯 & 40.53 & 38.48 & 48.84 & 59.20 & 48.75 & +20.3\% & 启动后回踩确认 \\
江波龙 & 587.60 & 520.00 & 900.00 & 1100.00 & 802.30 & +36.5\% & 高弹性回撤配置 \\
恒瑞医药 & 55.75 & 64.00 & 78.00 & 90.00 & 74.91 & +34.4\% & 趋势回撤核心 \\
工业富联 & 66.27 & 62.40 & 84.60 & 113.75 & 82.29 & +24.2\% & 业绩兑现回撤核心 \\
\bottomrule
\end{tabularx}
\sourcenote{AStock当前价多锚估值。静默篮子采用70\%基本面、20\%市场、10\%外部目标；已启动篮子采用60--65\%基本面、25--30\%市场、10\%外部目标，以反映更高拥挤风险。}
\end{exhibitbox}

\section{外部锚与内部冷判断}

\begin{exhibitbox}[Exhibit 23：券商预测与目标锚]
\begin{tabularx}{\textwidth}{L{2.2cm}L{2.7cm}L{2.5cm}L{2.2cm}X}
\toprule
\textbf{公司} & \textbf{外部预测} & \textbf{外部目标} & \textbf{方法} & \textbf{AStock处理} \\
\midrule
渝农商行 & 开源：2026E净利128.8亿元、EPS 1.13 & 华创：8.82元 & 0.72x PB & 目标和分母分开复核，Street权重10\% \\
徐工机械 & 多家EPS 0.67--0.72 & 华创：12.50元 & 18x PE & 取EPS 0.71，基准倍数降至16x \\
沪农商行 & 国信：净利124亿元、EPS 1.29 & 国泰海通：10.73元 & 0.75x PB & 基准仅0.70x，承认盈利增速偏低 \\
中兴通讯 & 当前共识：净利71亿元、EPS 1.48 & 群益：60元 & PE & 使用更保守分母，最终目标48.75元 \\
江波龙 & H1预告后：净利200亿元、EPS 47.28 & 大摩：673元 & 周期PE & AStock基准900元，但市场/周期权重更高 \\
恒瑞医药 & 华泰：净利94亿元、EPS 1.42 & 华泰87.14/高盛79.20 & SOTP & 最终74.91元，低于两家外部锚 \\
工业富联 & 华泰：净利560.1亿元、EPS 2.82 & 华泰：93元 & 33x PE & 基准仅30x；华创645.97亿元只作牛市盈利交叉验证 \\
\bottomrule
\end{tabularx}
\end{exhibitbox}

最终目标不是券商目标的平均数。中兴采用比旧券商更保守的利润分母；江波龙AStock基本面价值高于大摩目标，但市场锚权重提高至30\%；恒瑞最终目标低于华泰和高盛，避免把全部管线与BD乐观情景一次性计入；工业富联对华泰33倍PE下调到30倍，华创因未披露点目标而不进入Street目标权重。

\section{风险收益与组合角色}

\begin{itemize}
  \item \textbf{静默篮子稳健优先}：渝农商行；赔率优先：徐工机械；收益优先：沪农商行。
  \item \textbf{启动篮子弹性优先}：江波龙；增长质量优先：恒瑞医药；规模兑现优先：工业富联；均衡修复优先：中兴通讯。
\end{itemize}

如果构建研究观察组合，可先在两篮子之间各分配50\%；静默篮子内部按40\%渝农商行、35\%徐工机械、25\%沪农商行，启动篮子内部按30\%恒瑞、25\%江波龙、25\%工业富联、20\%中兴。该比例只表达研究信念，不是交易指令。
